\partstart{Relational Objects and Relational Complexity}
\section{Relation}\label{sec:relation}
\setcounter{theorem}{0}
\begin{definition}[Relation]
A relation is a locally available adjacency between distinguishable states or transducers through which an admissible event could occur.
\end{definition}

Carrier adjacency gives availability, while a domain $\mathcal D_{st}$ specifies which inputs actually permit a handoff. Objects are treated as persistent relational organizations. Effective laws are stable constraints on admitted compositions or their readouts; the admission rule must be fixed independently of the physical result being sought.

\begingroup
\renewcommand{\thesection}{4A}
\section{Registered Objects and Scalar Complexity}\label{sec:complexity}
\status{working definitions and requirements; selection of a general native functional remains a construction target}
\begin{definition}[Registered relational object]
A registered relational object $R$ is a specified relational structure together with its registration interface and the distinctions retained by that interface. Its description must declare its domain, constituent relations, relevant context, and the equivalence $\simeq_{\rm reg}$ under which two descriptions represent the same registered object.
\end{definition}
The object may be a structured carrier state, an admitted event with its receipt, or a registered history. These are different types and are not silently interchanged. A projected object can omit distinctions that a complete transduction must retain elsewhere. In particular, $R$ is not the reconstruction map $\mathsf R_e$, the receipt space $\mathcal R_e$, or an abstract based-return operator merely renamed.

\begin{definition}[Relational complexity]
Relative to the declared object domain and registration, a relational-complexity functional is a map
\begin{equation}
\mathcal C:\mathfrak R/\simeq_{\rm reg}\longrightarrow\R_{\geq0}.
\tag{RC1}\label{eq:complexitydomain}
\end{equation}
Its intended meaning is the relational complexity of the thing: how much relational organization is present in the registered object. Its construction must specify what it counts or measures, its normalization, and its behavior under admitted transformations. Equivalent registered presentations have the same value.
\end{definition}
This definition states a type and invariance requirement, not a unique formula. A native realization must be derived from the admitted structure rather than fitted to a desired physical interpretation. Two objects can have the same scalar without being the same relational object; a scalar need not encode every aspect of organization.

Raw distinct-vertex cardinality is only one possible support statistic, and an operator has no vertex set until a support convention is supplied. History length, net winding, continuation growth, receipt capacity, and entropy are likewise distinct quantities. No universal identification of these with $\mathcal C$ is made. Redundant copies of a description do not increase the complexity of the same registered object; an enlarged object with genuinely additional registered relations is a different case.

\begin{definition}[Scalar encoding]
For a fixed dimensionless encoding base $G>1$, the associated multiplicative scalar is
\begin{equation}
S(R)=G^{\mathcal C(R)},\qquad \log_G S(R)=\mathcal C(R).
\tag{RC2}\label{eq:scalarencoding}
\end{equation}
The logarithmic coordinate $\mathcal C$ and the positive scalar $S$ contain the same scalar information. The exponential is optional notation, not an additional physical law.
\end{definition}
Here $G$ is neither a graph such as $G_{60}$ nor Newton's $G_N$. Likewise $S(R)$ is not the dimensionful mechanical action $S[\gamma]$ of Section~\ref{sec:hbar}. A change of encoding base changes the numerical representation, not the underlying relational object. With the present nonnegative convention $S\geq1$; negative complexity increments are nevertheless possible.

\subsection*{Complexity and information have different roles}
Conservation of information constrains the complete input-output map. It does not require every structural readout to be constant. An invertible reorganization may change a chosen complexity functional while preserving all input distinctions in the joint output. Conversely, a constant scalar does not prove preservation: a many-to-one map can preserve a coarse count while losing distinctions.

Phase and Mode describe organization or orientation, not an extra stock of information to be added to $\mathcal C$. No identity of the form $\mathcal I=\mathcal C+\mathcal I_{\rm phase}+\mathcal I_{\rm Mode}$ is assumed. The joint identities in Section~\ref{sec:information} already account for correlations, and remain the quantitative statements available when the required classical distribution or quantum state is specified.

\subsection*{Conserved information and invariant readouts}
Complete recovery can coexist with a changing $\mathcal C$; an unchanged $\mathcal C$ can coexist with changing orientation; and a constant scalar alone cannot certify recovery. The manuscript therefore distinguishes the general information principle from an additional invariant readout. The supplied quadratic $I_{\mathsf G}(X)=X^\dagger\mathsf G X$ in Section~\ref{sec:conservativedynamics} is one such readout. No identification with $\mathcal C(R)$, $G^{\mathcal C(R)}$, entropy, or physical energy is assumed without a separate construction.
\endgroup
\setcounter{section}{4}

\begingroup
\renewcommand{\thesection}{4B}
\section{Selected Reference and Situated Registration}\label{sec:executingreference}
\status{operational definitions; signed-face incidence reconstructed from supplied native data}
An account is registered somewhere in the declared apparatus. The reader reconstructing an experiment and the apparatus executing it each use a particular presentation and retain a particular record. Neither must occupy a position privileged by the laws. What is required is an identified interface at which an input can be compared with a retained consequence. The distinction is between a selected reference and a physically preferred frame.

\begin{definition}[Executing reference]\label{def:executingreference}
An executing reference is a selected presentation at a specified registration interface, together with the retained context and any orientation data required by the intended readout. Write
\begin{equation}
\rho_*=(\mathsf{Reg}_*,p_*,m_*,\vartheta_*),
\tag{RF1}\label{eq:executingreference}
\end{equation}
where $\mathsf{Reg}_*$ specifies the interface, $p_*$ the selected presentation, $m_*$ its accessible retained context, and $\vartheta_*$ the required orientation convention, when one is needed. The admitted domain and rules for changing these data are part of the specification.
\end{definition}
A selected field address is not automatically a complete frame or an orientation above that address. If a readout requires additional data, selecting the address alone does not complete it. Nor does naming an executor provide a new event law. The complete update remains $\mathsf T_e:(x,m)\mapsto(p,m')$ as in equation~(7).

The motivating statement is \emph{``History accumulates here. I am 1.''} Its mathematical use is a normalization of a selected reference, not an identity involving the operator $I$ or a claim that the executor has unit complexity. On a two-member orientation torsor, the typed version will be $\sigma_{t_*}(t_*)=+1$. This sets neither $\mathcal C(R)=1$ nor $\Delta\mathcal C=1$, supplies no clock calibration, and does not rederive $\cth=1$.

Here \emph{subjective} means situated at a specified record-bearing reference. It does not mean erroneous or unsigned. An \emph{objective comparison}, in the operational sense used here, is a comparison whose stated content is consistent under the admitted changes of reference. It need not be a view from nowhere or an inventory held completely by any one participant. The invariance or covariance required of that comparison must be proved for its actual interface.

\subsection*{Why incidence belongs in the presentation}
A state and its occurrence at a boundary are different types. The supplied native signed-face system has twenty-four five-state blocks, arranged in six carriers with two positive and two negative blocks on each carrier. Each sign section partitions the sixty states. Consequently the native boundary domain is
\begin{equation}
\mathcal I_F=\{(u,B):u\in B\},\qquad |\mathcal I_F|=120.
\tag{RF2}\label{eq:faceincidence}
\end{equation}
Each state $u$ has two such presentations, one of each sign. The sign is therefore not a function of $u$ alone: it belongs to its incidence with the selected block. The signs here are the supplied boundary/event-side signs, not yet the orientation signs of a Pauli analyzer. The original signed-block record and Program~2's projective-line record are included without modification \citep{cave2026nativepacket}.

The four blocks on a carrier admit exactly two partitions into a pair of positive--negative pairs: they are the two bijections between its two positive and two negative blocks. Denote their unordered set by $\mathcal P(f)$. The phase-partition-decorated incidence is
\begin{equation}
\widetilde{\mathcal I}_F=
\{(u,B,P):(u,B)\in\mathcal I_F,\ P\in\mathcal P(f(B))\},
\qquad |\widetilde{\mathcal I}_F|=240.
\tag{RF3}\label{eq:decoratedfaceincidence}
\end{equation}
Every bare incidence remains compatible with both partitions. Thus incidence is necessary context for this presentation, but bare incidence does not select its independent phase. The twelve partitions reconstructed here agree exactly with the supplied Program~2 Audit018 partition table, without identifying its selected $\epsilon$ labels by record order.

For a native permutation $g$, the set-theoretic action is
\[
(u,B,P)\longmapsto(gu,gB,gP).
\]
The included replay checks all 480 cached permutations, closure on signed blocks and partitions, and equivariance of forgetting $P$. Composition is checked on a generating set against every group element. It does not infer an independent phase selector, an analyzer direction, or a chronological source from these counts \citep{cave2026draft6replays}.

\subsection*{A reference supplies coordinates, not missing information}
Two situated records can be reconciled only through an admitted comparison interface. Knowing that two presentations exist is not yet knowing how their orientations are related. Section~\ref{sec:incidencereconciliation} makes the distinction exact: a selected local reference gives sign coordinates, while an incidence-dependent comparison must obey its own covariance law.

This requirement also protects the scalar-relay and chronology boundaries. A receiver still needs the data required by Proposition~11A.3; calling its account complete does not make unavailable information local. A record can grow while its finite controller ignores it. Naming that record ``my history'' does not create the nonperiodic continuation mechanism required in Section~\ref{sec:frequency}.
\endgroup
\setcounter{section}{4}


\section{Continuation}
\setcounter{theorem}{0}
\begin{definition}[Continuation]
A continuation is a lawful possible next relation available from the present relational and registration context.
\end{definition}

For a complete context $z$, the set $\Adm(z)$ contains its admitted elementary continuations. When a fixed matrix $A_{\rm cont}$ counts labelled transitions or their declared multiplicities, the number of length-$L$ words from $i$ is
% Original equation number(s) 8; expression unchanged.
\setcounter{equation}{7}
\begin{equation}
N_L(i)=\bigl(A_{\rm cont}^{L}\mathbf 1\bigr)_i.
\end{equation}

An associated growth rate is
% Original equation number(s) 9; expression unchanged.
\setcounter{equation}{8}
\begin{equation}
h_{\rm cont}(i)=\limsup_{L\to\infty}\frac{1}{L}\log N_L(i),
\end{equation}

This measures growth of the admitted language. A regular $k$-fold cover multiplies the count of lifts of each base path by the constant $k$ and therefore leaves this exponential rate unchanged. Cover multiplicity alone supplies no additional entropy rate.

\section{Informative Action}\label{sec:action}
\setcounter{theorem}{0}
\begin{definition}[Informative action]
Informative action is an admissible realized continuation that creates or changes a registered relational distinction.
\end{definition}

Informativeness is relative to a stated registration or future-response family. A lifted return with unchanged projected endpoint can be informative through its deck residue. Primitive informative action $\mathcal A_e$ is distinct from a dimensionful mechanical functional $S[\gamma]$; their correspondence is addressed in Section~\ref{sec:hbar}.

\begingroup
\renewcommand{\thesection}{6A}
\section{Informative Coupling and Information Current}\label{sec:coupling}
\status{operational definitions; finite illustration on the declared relational kernel}
Informative action is a realized continuation. Coupling instead specifies a dependence that a continuation can use. The distinction is important when an extra cover coordinate is available but has no demonstrated effect on the future.

\begin{definition}[Informative coupling]
Relative to a declared context $z$ and retained variables $r,s$, a law is informatively coupled to $r$ beyond $s$ if changing $r$ while fixing $z,s$ changes an admitted continuation set or a registered response. In a deterministic description a witness is
\begin{equation}
\Adm(z,r,s)\ne\Adm(z,r',s)
\quad\text{or}\quad
\mathcal O_{z,r,s}(w)\ne\mathcal O_{z,r',s}(w)
\tag{IC1}\label{eq:informativecoupling}
\end{equation}
for a common admitted test $w$. A stochastic formulation compares response distributions under the same declared intervention and preparation conditions.
\end{definition}
This definition concerns operational dependence, not mere coexistence or correlation of labels. In the kernel example $F_0$, changing $x_1$ changes the allowed target $y_2$, providing an exact witness. The kernel relation alone does not choose a target, retain all complementary input information or supply an adjacent field handoff.

\begin{definition}[Information current]
An information current is transport of an operationally distinguishable payload through admitted field handoffs, with specified encoding, continuation and downstream decoding. Each participating handoff must preserve a dependence of its downstream response on the incoming distinction, within the declared context and complete-output accounting.
\end{definition}
Under these definitions an actual information current uses informative coupling, while coupling can remain local without a current. ``Current'' is not merely a nonzero parity label, a static correlation or electric charge current by renaming. Quantitative information flux would additionally require a measure, an interface and an event or time normalization.

For the cover variable $\epsilon_K$, informative coupling remains a separate test:
\[
\mathcal O_{z,0}(w)\stackrel{?}{\ne}\mathcal O_{z,1}(w)
\quad\text{or}\quad
\Adm(z,0)\stackrel{?}{\ne}\Adm(z,1).
\]
Here $z$ must explicitly state which history and other variables are held fixed. The canonical double gives the two sheet values and their walk-parity update; it does not establish either inequality. The realized effect of a sheet-sensitive law, rather than the existence of a second sheet, would make that extra distinction active in the transducer.
\endgroup
\setcounter{section}{6}

\begingroup
\renewcommand{\thesection}{6B}
\section{Conservative Informative Action}\label{sec:conservativeaction}
\status{definition within the existing admissible-transduction framework}
Coupling specifies a dependence; action realizes a continuation; current additionally requires an encoded distinction, admitted field handoffs, and downstream decoding. Conservation is a property of the complete update, not a fourth name for any of these objects.

\begin{definition}[Conservative informative action]\label{def:conservativeaction}
Fix a complete registered state domain, its admitted recoverable transductions, a readout family that witnesses informativeness, and a declared invariant $I$. An informative action is conservative relative to $I$ when its realized transduction $\mathsf T$ satisfies $I(\mathsf Tz)=I(z)$ on its admitted domain.
\end{definition}
The invariant may be one readout of a much richer object. It need not equal $\mathcal C$, and invariance alone neither certifies recovery nor supplies an informative effect. In a faithful linear realization a useful special case is $I_{\mathsf G}(X)=X^\dagger\mathsf G X$. Section~\ref{sec:conservativedynamics} proves the associated skew-generator condition and its limits. ``Informational form'' is justified only when the representation and its relation to registered distinctions have also been declared.

Gyroscopic orientation and source-free electromagnetic evolution provide two conventional comparison models. They show how an invariant can constrain a changing state; they do not establish that every admissible QR action is linear, metric-preserving, or physical. Likewise an internal conservative update need not be an adjacent handoff and does not by itself increment propagation depth.
\endgroup
\setcounter{section}{6}


\partstart{Information-Preserving Transduction}
\section{Admissible Transduction}\label{sec:transduction}
\setcounter{theorem}{0}
\begin{definition}[Admissible transduction]
An admissible transduction is a local transformation of a complete incoming relational state that satisfies the declared native constraints and preserves the information required for a reversible admissible handoff.
\end{definition}

In abbreviated classical notation,
% Original equation number(s) 10; expression unchanged.
\setcounter{equation}{9}
\begin{equation}
\mathsf T_e:\mathcal D_e\longrightarrow\mathcal P_e\times\mathcal R_e,
\qquad x\longmapsto(p,r).
\label{eq:primitiveT}
\end{equation}

with reconstruction on the admitted image:
% Original equation number(s) 11; expression unchanged.
\setcounter{equation}{10}
\begin{equation}
\mathsf R_e\mathsf T_e=\id_{\mathcal D_e}.
\label{eq:leftinverse}
\end{equation}

Surjectivity onto all abstract output pairs is unnecessary. If the input contains old memory $m$, the reconstructed object is $(x,m)$, not just the newly arriving payload.

\setcounter{theorem}{1}
\begin{example}[Propagation with a nontrivial receipt]
For $u,v\in\F_2$, let $p=u+v$ and $r=v$. Then
\[
\mathsf T(u,v)=(u+v,v),\qquad \mathsf R(p,r)=(p+r,r).
\]
The propagated bit alone cannot distinguish $(0,0)$ from $(1,1)$, but the receipt restores the distinction. No information has been copied: two input bits have been reorganized into two output bits.
\end{example}

The two-bit example illustrates the accounting form. Carrier-specific operations require their own admission rule. In a coherent quantum representation the corresponding complete map is an isometry,
% Original equation number(s) 12; expression unchanged.
\setcounter{equation}{11}
\begin{equation}
V_e:\mathcal H_X\longrightarrow\mathcal H_P\otimes\mathcal H_R,
\qquad V_e^\dagger V_e=I.
\end{equation}

with recovery by $V_e^\dagger$ on its image. The joint output can be entangled; tracing out a receipt gives a channel rather than a generally reversible local map \citep{watrous2018}.

\section{The Propagated Component}
The propagated component is eligible to acquire a new adjacent field address under the admitted handoff rule. An internal carrier coordinate and a field address have different roles: a Thalion can have sixty internal states at one field site. The address graph $\mathcal F$ must be specified.

A transported distinction remains identifiable through declared encoders and decoders even when its coordinates change. Successive handoffs are new occurrences carrying the same relational content or event type, rather than one completed occurrence moving between addresses.

\section{Trace, Receipt, and Relational Residue}\label{sec:receipts}
A trace is a retained relational consequence. It functions as a receipt when, jointly with the propagated component and prior context, it accounts for the handoff. Attestation and reconstruction are different interfaces: a short record can attest an occurrence without reconstructing its complete microscopic input.
% Original equation number(s) 13; expression unchanged.
\setcounter{equation}{12}
\begin{equation}
\mathsf T_e\longrightarrow\delta_e\longrightarrow\Delta[\Gamma]\longrightarrow\tau[\Gamma]
\end{equation}

These arrows denote composition or readout maps. A local distinction $\delta_e$, a retained history $\Delta[\Gamma]$, and a temporal scalar are different objects. A finite deck residue stores only a shadow of an arbitrarily long history.

When the propagated map is already injective, a constant receipt can suffice. For quantum payloads, the receipt must respect the no-cloning and undisturbed-payload restrictions \citep{wootterszurek1982}; the precise proposition appears in Section~\ref{sec:quantum}. The term \emph{Relational Spin Residue} is reserved for a separately established spin-related central or projective residue. An arbitrary binary receipt does not acquire that interpretation by its order alone.

\section{The QR Event}\label{sec:event}
\setcounter{theorem}{0}
\begin{definition}[QR propagation event]
A QR event, in the propagation architecture studied here, is one completed, receipted, admissible relational handoff: a local transduction followed by an admitted adjacent transmission, with the information required to reconstruct the complete handoff retained.
\end{definition}

Internal comparison and registration can be stages of a propagation event; only the completed admitted adjacent handoff counts as one primitive propagation step.

An event is complete when the transmitted component reaches the next admitted address and the complete output preserves the local accounting. A complexity change is a readout of this transformation, not a replacement for its admission, handoff, or receipt requirements. Registration is an interaction represented within the apparatus. Its retained consequence need not already be available at another address. Historical placement is supplied by composition or causal dependency.

\emph{Event primacy} organizes the theory around these completed relational occurrences. States describe their input, output, and intermediate conditions. Projected snapshots alone need not identify the admitted history that produced them.
